% Source: Wald GR, physical PDF pp. 101--103
%         (printed pp. 88--90).
% Coverage: Chapter 4 problems 1--9.
% Figures/tables/footnotes: no figures, tables, or footnotes.
% Status: source-reviewed and compile-clean.
% Uncertainties: none.

\chapterproblems

\noindent\textbf{1.}\label{prob:4.1}\quad Show that Maxwell's
equation~\eqref{eq:4.3.12} implies strict charge conservation,
\(\nabla_aj^a=0\).

\medskip
\noindent\textbf{2.}\label{prob:4.2}\quad
\textbf{(a)} Let \(\alpha\in\Omega^p(M)\) be a \(p\)-form on an
\(n\)-dimensional oriented manifold with metric \(g_{ab}\), i.e.,
\(\alpha_{a_1\dots a_p}\) is a totally antisymmetric tensor field (see appendix B).
We define the dual, \(\star\alpha\), of \(\alpha\) by
\begin{equation*}
  (\star\alpha)_{b_1\dots b_{n-p}}
    =\frac{1}{p!}\alpha^{a_1\dots a_p}
      \epsilon_{a_1\dots a_pb_1\dots b_{n-p}} ,
\end{equation*}
where \(\epsilon_{a_1\dots a_n}\) is the natural volume element on \(M\), i.e.,
the totally antisymmetric tensor field determined up to sign by
equation~\eqref{eq:B.2.9}. Show that
\(\star\star\alpha=(-1)^{s+p(n-p)}\alpha\), where \(s\) is the number of minuses
occurring in the signature of \(g_{ab}\).

\noindent\textbf{(b)} Show that in differential forms notation (see appendix B),
Maxwell's equations~\eqref{eq:4.3.12} and \eqref{eq:4.3.13} can be written as
\begin{align*}
  \dd\star F&=4\pi\star j ,\\
  \dd F&=0 .
\end{align*}
Note that if we apply Stokes's theorem (see appendix B) to the first equation, we
obtain \(\int_\Sigma\star j=(1/4\pi)\int_S\star F\), where \(\Sigma\) is a
three-dimensional hypersurface with two-dimensional boundary \(S\). But
\(-\int_\Sigma\star j=-\int_\Sigma j^at_a\dd\Sigma\) is just the total electric
charge \(e\) in the volume \(\Sigma\), where \(t^a\) is the unit normal to
\(\Sigma\), and \(-\int_S\star F=\int_S E^an_a\dd A\) is just the integral of the
normal component of \(E_a=F_{ab}t^b\) on \(S\). Thus, Gauss's law of
electromagnetism continues to hold in curved spacetime.

\noindent\textbf{(c)} Define for each \(\beta\in[0,2\pi]\) the tensor field
\(\widetilde F_{ab}=F_{ab}\cos\beta+(\star F)_{ab}\sin\beta\). We call
\(\widetilde F_{ab}\) a \emph{duality rotation} of \(F_{ab}\) by ``angle''
\(\beta\). It follows immediately from part (b) that if \(F_{ab}\) satisfies the
source-free Maxwell's equations (\(j^a=0\)), then so does \(\widetilde F_{ab}\).
Show that the stress-energy, \(T_{ab}\), of the solution \(\widetilde F_{ab}\) is the
same as that of \(F_{ab}\).

\medskip
\noindent\textbf{3.}\label{prob:4.3}\quad
\textbf{(a)} Derive equation~\eqref{eq:4.4.24}.

\noindent\textbf{(b)} Show that the ``gravitational electric and magnetic fields''
\(\mathbf{E}\) and \(\mathbf{B}\) inside a spherical shell of mass \(M\) and radius \(R\)
(with \(M\ll R\)) slowly rotating with angular velocity \(\boldsymbol{\omega}\) are
\begin{equation*}
  \mathbf{E}=0 ,\qquad \mathbf{B}=\frac{2M}{3R}\boldsymbol{\omega} .
\end{equation*}

\noindent\textbf{(c)} An observer at rest at the center of the shell of part (b)
parallelly propagates along his (geodesic) world line a vector \(S^a\) with
\(S^au_a=0\), where \(u^a\) is the tangent to his world line. Show that the inertial
components, \(\mathbf{S}\), precess according to
\(\dd\mathbf{S}/\dd t=\boldsymbol{\Omega}\mathbin{\times}\mathbf{S}\), where
\(\boldsymbol{\Omega}=2\mathbf{B}=\tfrac43(M/R)\boldsymbol{\omega}\). This effect, first analyzed by
Thirring and Lense (1918) and discussed further by Brill and Cohen (1966), may be
interpreted as a ``dragging of inertial frames'' caused by the rotating shell. At the
center of the shell the local standard of ``nonrotating,'' defined by parallel
propagation along a geodesic, is changed from what it would be without the shell, in
a manner in accord with Mach's principle.

\medskip
\noindent\textbf{4.}\label{prob:4.4}\quad Starting with
equation~\eqref{eq:3.4.5} for \(R_{ab}\), derive the formula,
equation~\eqref{eq:4.4.51}, for \(R^{(2)}_{ab}\) by substituting
\(\eta_{ab}+h_{ab}\) for \(g_{ab}\) and keeping precisely the terms quadratic in
\(h_{ab}\).

\medskip
\noindent\textbf{5.}\label{prob:4.5}\quad Let \(T_{ab}\) be a symmetric,
conserved tensor field (i.e., \(T_{ab}=T_{ba}\), \(\partial^aT_{ab}=0\)) in
Minkowski spacetime. Show that there exists a tensor field \(U_{acbd}\) with the
symmetries \(U_{acbd}=U_{[ac]bd}=U_{ac[bd]}=U_{bdac}\) such that
\(T_{ab}=\partial^c\partial^dU_{acbd}\). (Hint: For any vector field \(v^a\) in
Minkowski spacetime satisfying \(\partial_av^a=0\) there exists a tensor field
\(s^{ab}=-s^{ba}\) such that \(v^a=\partial_bs^{ab}\). [This follows from applying
the converse of the Poincar\'e lemma (see the end of section B.1 in appendix B) to
the 3-form \(\epsilon_{abcd}v^d\).] Use this fact to show that
\(T_{ab}=\partial^cW_{cab}\) where \(W_{cab}=W_{[ca]b}\). Then use the fact that
\(\partial^cW_{c[ab]}=0\) to derive the desired result.)

\medskip
\noindent\textbf{6.}\label{prob:4.6}\quad As discussed in the text, in general
relativity no meaningful expression is known for the local stress-energy of the
gravitational field. However, a four-index tensor \(T_{abcd}\) can be constructed
out of the curvature in a manner closely analogous to the way in which the stress
tensor of the electromagnetic field is constructed out of \(F_{ab}\)
(equation~\eqref{eq:4.2.27}). We define the \emph{Bel--Robinson tensor} in terms
of the Weyl tensor by
\begin{align*}
  T_{abcd}
    &=C_{aecf}C_b{}^e{}_d{}^f
      +\frac14\epsilon_{ae}{}^{hi}\epsilon_b{}^{ejk}
        C_{hicf}C_{jkd}{}^f\\
    &=C_{aecf}C_b{}^e{}_d{}^f
      -\frac32g_{a[b}C_{jk]cf}C^{jk}{}_d{}^f ,
\end{align*}
where \(\epsilon_{abcd}\) is defined in appendix B and
equation~\eqref{eq:B.2.13} was used. It follows that
\(T_{abcd}=T_{(abcd)}\). (This is established most easily from the spinor
decomposition of the Weyl tensor given in chapter 13.)

\noindent\textbf{(a)} Show that \(T^a{}_{acd}=0\).

\noindent\textbf{(b)} Using the Bianchi identity~\eqref{eq:3.2.16}, show that in
vacuum, \(R_{ab}=0\), we have \(\nabla^aT_{abcd}=0\).

\medskip
\noindent\textbf{7.}\label{prob:4.7}\quad
\textbf{(a)} Show that the total energy \(E\), equation~\eqref{eq:4.4.55}, is time
independent, i.e., the value of \(E\) is unchanged if the integral is performed over
a time translate, \(\Sigma'\), of \(\Sigma\).

\noindent\textbf{(b)} Let \(\xi_a\) be a gauge transformation which vanishes
outside a bounded region of space. Show that
\(E[h_{ab}]=E[h_{ab}+2\partial_{(a}\xi_{b)}]\) by comparing them with
\(E[h_{ab}+2\partial_{(a}\xi'_{b)}]\) where \(\xi'_a\) is a new gauge
transformation which agrees with \(\xi_a\) in a neighborhood of the hyperplane
\(\Sigma\) but vanishes in a neighborhood of another hyperplane \(\Sigma'\).

\medskip
\noindent\textbf{8.}\label{prob:4.8}\quad Two point masses of mass \(M\) are
attached to the ends of a spring of spring constant \(K\). The spring is set into
oscillation. In the quadrupole approximation, equation~\eqref{eq:4.4.58}, what
fraction of the energy of oscillation of the spring is radiated away during one cycle
of oscillation?

\medskip
\noindent\textbf{9.}\label{prob:4.9}\quad A binary star system consists of two
stars of mass \(M\) and of negligible size in a nearly Newtonian circular orbit of
radius \(R\) around each other. Assuming the validity of
equation~\eqref{eq:4.4.58} for this system, calculate the rate of increase of the
orbital frequency due to emission of gravitational radiation.
